\chapter{Introdução e Contexto do Domínio de Dados}
\label{cap_introducao_dados}

O desenvolvimento do Sistema Inteligente de Gestão de Eventos Adversos baseado na metodologia \textit{Global Trigger Tool} (SIGEA-GTT) insere-se em um contexto interdisciplinar crítico, integrando a Ciência da Computação (Departamento de Computação -- DCOMP/UFS) e as Ciências da Saúde (Departamento de Enfermagem -- UFS). A plataforma foi concebida para atender à necessidade de formação sistemática de profissionais e acadêmicos na identificação, rastreamento e análise de eventos adversos (EA) em prontuários hospitalares, empregando o padrão estabelecido pelo \textit{Institute for Healthcare Improvement} (IHI) \cite{ihi2009}.

Em sistemas de informação biomédicos e educacionais dessa natureza, a modelagem da camada de persistência constitui a base fundamental para a confiabilidade operacional, a integridade clínica dos dados, o suporte a auditorias concorrentes e a estrita conformidade legal com a Lei Geral de Proteção de Dados Pessoais (LGPD -- Lei nº 13.709/2018) \cite{lgpd2018}. Este documento detalha formalmente o projeto do banco de dados do SIGEA-GTT, articulando desde a modelagem conceitual abstrata até a implementação físico-relacional no SGBD PostgreSQL 16 LTS \cite{postgresql2024}.

% ===========================================================================
\section{A Tríplice Fronteira de Dados do Sistema}
\label{sec_triplice_fronteira}
% ===========================================================================

O domínio de dados do SIGEA-GTT é caracterizado por uma tríplice fronteira de requisitos informacionais, exigindo uma arquitetura de banco de dados capaz de harmonizar o rigor transacional relacional com a flexibilidade de documentos clínicos complexos:

\begin{enumerate}
    \item \textbf{Fronteira de Segurança, Identidade e Acesso (RBAC)}:
    \begin{itemize}
        \item Gestão de contas de usuários restritas ao domínio institucional \texttt{@academico.ufs.br};
        \item Controle de Acesso Baseado em Papéis (\textit{Role-Based Access Control} -- RBAC) restrito a três perfis: Administrador (\texttt{ADMIN}), Docente (\texttt{PROFESSOR}) e Discente (\texttt{STUDENT});
        \item Regra de negócio de matrícula: o identificador institucional acadêmico (matrícula) é de preenchimento obrigatório apenas para o discente, sendo estritamente nulo para gestão e docência.
    \end{itemize}

    \item \textbf{Fronteira de Taxonomia Clínica IHI-GTT e Gravidades NCC MERP}:
    \begin{itemize}
        \item Catálogo canônico dos seis módulos analíticos do IHI: Cuidados Gerais (\textit{General Care}), Cirúrgico (\textit{Surgical}), Medicamentos (\textit{Medication}), Terapia Intensiva (\textit{Intensive Care}), Perinatal (\textit{Perinatal}) e Emergência (\textit{Emergency Department});
        \item Cadastro dos 53 gatilhos clínicos padronizados vinculados hierarquicamente aos seus módulos;
        \item Classificação de gravidade de dano estabelecida pelo conselho NCC MERP \cite{nccmerp2001}, segregando incidentes sem dano (Categorias A a D) e eventos com dano real (Categorias E a I).
    \end{itemize}

    \item \textbf{Fronteira Pedagógica, Prontuários Simulados e Ferramentas da Qualidade}:
    \begin{itemize}
        \item Gestão de turmas acadêmicas identificadas pela tripla canônica \textit{"Nome da Matéria -- Código da Turma -- Período Letivo"} e associação de discentes matriculados;
        \item Criação de atividades avaliativas vinculadas a prontuários eletrônicos simulados de alta fidelidade;
        \item Submissões de auditoria discentes contendo rastreamento temporal sob cronômetro de 20 minutos, gatilhos detectados com justificativa clínica, gravidade atribuída e preenchimento de ferramentas da qualidade (Diagrama de Ishikawa, Matriz GUT, Plano 5W2H e Ciclo PDCA), além de nota e parecer docente.
    \end{itemize}
\end{enumerate}

% ===========================================================================
\section{Padrões de Modelagem e Paradigma Adotado}
\label{sec_padroes_adotados}
% ===========================================================================

Para garantir clareza pedagógica e rigor de engenharia de software, este projeto adota uma progressão clássica de modelagem em três níveis de abstração \cite{elmasri2016, heuser2009}:

\begin{enumerate}
    \item \textbf{Nível Conceitual (MER)}: Formalizado com base na teoria de Peter Chen (1976) \cite{chen1976}, identificando entidades fortes, fracas e associativas, atributos identificadores e relacionamentos com cardinalidades explícitas em diagrama nativo \textit{TikZ};
    \item \textbf{Nível Lógico (DER)}: Formalizado segundo o modelo relacional de Edgar F. Codd (1970) \cite{codd1970} e a notação Pé de Galinha (\textit{Crow's Foot}) de James Martin (1989) \cite{martin1989}, compreendendo a definição de chaves primárias, chaves estrangeiras, tipos de dados e dicionário de dados relacional completo;
    \item \textbf{Nível Físico}: Implementação no SGBD PostgreSQL 16 LTS \cite{postgresql2024}, fundamentada em uma abordagem híbrida relacional-documental (\textit{Hybrid Relational-JSONB Design}), combinando integridade referencial estrita e normalização até a BCNF com colunas binárias \texttt{JSONB} indexadas via GIN para casos clínicos simulados flexíveis.
\end{enumerate}

% ===========================================================================
\section{Organização do Documento}
\label{sec_organizacao_doc}
% ===========================================================================

O restante deste documento está estruturado conforme segue:
o Capítulo \ref{cap_modelo_conceitual} detalha o Modelo Entidade-Relacionamento Conceitual (MER) e seu diagrama canônico;
o Capítulo \ref{cap_normalizacao_regras} apresenta a demonstração formal de normalização (1FN a BCNF), a justificativa técnica para uso de \texttt{JSONB} e as regras de integridade de domínio;
o Capítulo \ref{cap_modelo_logico} exibe o Diagrama Entidade-Relacionamento Lógico (DER) e o Dicionário de Dados exaustivo das oito tabelas físicas;
o Capítulo \ref{cap_esquemas_json} detalha formalmente a especificação dos esquemas \textit{JSON Schema} das estruturas semiestruturadas;
o Capítulo \ref{cap_modelo_fisico} analisa o modelo físico, estratégias de indexação B-Tree/GIN e integridade referencial;
o Capítulo \ref{cap_matriz_rastreabilidade} expõe a matriz de rastreabilidade entre requisitos e modelo de dados;
e o Capítulo \ref{cap_consideracoes_finais} sintetiza as conclusões, conformidade com a LGPD e perspectivas de governança de dados.
